% SAI / IEEE-style two-column conference format (IEEEtran.cls is not installed
% locally, so we emulate the SAI_PAPER_FORMAT template in the article class:
% two columns, US-letter, Roman-numeral section heads, "Abstract--"/"Keywords--"
% run-ins, and numeric references in citation order).
\documentclass[10pt,twocolumn,letterpaper]{article}

% --- packages (all confirmed present in the local TeX install) ---
\usepackage[utf8]{inputenc}
\usepackage[T1]{fontenc}
\usepackage{lmodern} % scalable fonts, required for microtype font expansion
\usepackage[letterpaper,top=0.75in,bottom=1in,left=0.625in,right=0.625in,columnsep=0.25in]{geometry}
\usepackage{amsmath,amssymb}
\usepackage{graphicx}
\usepackage{booktabs}
\usepackage{multirow}
\usepackage{caption}
\usepackage{microtype}
\usepackage{xcolor}
\usepackage[numbers,sort&compress]{natbib}
\usepackage{xurl} % allow long reference URLs to break across lines
\usepackage[hidelinks]{hyperref}

% IEEE/SAI section numbering: Roman sections (I, II, ...), lettered subsections.
\renewcommand{\thesection}{\Roman{section}}
\renewcommand{\thesubsection}{\Alph{subsection}}
\renewcommand{\thesubsubsection}{\arabic{subsubsection}}
\captionsetup{font=small,labelsep=period}

\newcommand{\gcp}{\textsc{gcp}}
\newcommand{\ata}{\textsc{a2a}}
\newcommand{\bigoh}{\mathcal{O}}

\title{\textbf{Read-First, Role-Gated Context Federation Beats\\
Point-to-Point Message-Passing for Multi-Agent LLM Systems:\\
A Fairness-Controlled Evaluation}}

\author{%
Adrian Issac Mamani Quispe\\
\textit{Universidad Nacional de San Agust\'in de Arequipa}\\
Arequipa, Peru\\
admamaniq@unsa.edu.pe
\and
Ryan Fabian Valdivia Segovia\\
\textit{Universidad Nacional de San Agust\'in de Arequipa}\\
Arequipa, Peru\\
rvaldiviase@unsa.edu.pe
}

\date{June 2026}

\begin{document}
\maketitle

\begin{abstract}
Multi-agent systems built on large language models (LLMs) overwhelmingly coordinate by
\emph{message-passing}: each agent sends task and context directly to named peers, as in
Google's Agent2Agent (A2A) protocol and the classical agent communication languages it
descends from. We argue that an alternative substrate---\emph{read-first, role-gated context
federation} (the Graph Context Protocol, \gcp{})---is structurally superior on three axes at
equal task success: \textbf{coupling}, \textbf{information leakage}, and \textbf{cost}. In
\gcp{}, an agent does not push context to peers; it \emph{reads} context from independently
owned knowledge nodes through a policy gate that authorizes the read by role and records a
provenance entry for every decision. We make the comparison \emph{fair} by holding everything
constant except the interop layer: both arms run the identical agent brain, the identical model,
and the identical task; only the per-peer tool factory and node hosting differ (an in-process
\gcp{} fetch handler versus real localhost A2A servers). On three scenarios---a marketplace
(coupling), a software organization (leakage), and a cross-owner supply chain---we find: (i) the
number of pairwise integration edges grows as $\bigoh(n^2)$ for \ata{} versus $\bigoh(n)$ for
\gcp{} ($2450$ vs.\ $50$ at $n{=}50$), with constant marginal integration effort for \gcp{};
(ii) provenance completeness is $1.0$ for \gcp{} and $0.0$ for \ata{} across all scenarios and
seeds, because read decisions are audited by construction whereas point-to-point messages leave
no access trail; and (iii) \gcp{} never costs more, and on the cross-owner task costs fewer
tokens with far lower variance ($707.6{\pm}3.5$ vs.\ $901.4{\pm}229.2$). Task success is equal
across arms. We position \gcp{} against five decades of coupling theory and access-control
research and against recent LLM-agent security work, and we release the harness, scenarios, and
seeds. The evaluation is preliminary---one free-tier model, five seeds---and we are explicit
about what that does and does not license.
\end{abstract}

\smallskip
\noindent\textbf{\textit{Keywords---}}\textit{context federation; multi-agent
systems; large language models; access control; least privilege; provenance;
software coupling; agent communication protocols}
\smallskip

\section{Introduction}

The dominant coordination pattern for LLM-based multi-agent systems is direct message-passing.
Frameworks such as AutoGen~\citep{wu2023autogen}, ChatDev~\citep{qian2024chatdev}, and
MetaGPT~\citep{hong2024metagpt} compose applications from agents that send and receive messages,
and the emerging Agent2Agent (A2A) protocol~\citep{a2a2025} standardizes exactly this: a client
agent advertises and delegates tasks to named remote agents over JSON-RPC. This paradigm is a
direct descendant of the classical agent communication languages KQML~\citep{finin1994} and
FIPA-ACL~\citep{fipa2002}, which already assumed pairwise connections and shared, pre-negotiated
protocols between every interacting pair.

Message-passing has a structural cost that the software-engineering literature has understood
for fifty years. Connecting $n$ components pairwise requires on the order of $n^2$ links;
introducing a shared intermediary that each component connects to once reduces this to $n$
links. This is the ``integration spaghetti'' that the Enterprise Integration
Patterns~\citep{hohpe2003} catalog addresses with message brokers and hubs, and it is the same
phenomenon that coupling theory~\citep{parnas1972,stevens1974,briand1999} formalizes: a system's
modifiability is governed by the number and strength of its inter-component dependencies.

We study an alternative substrate for multi-agent coordination, the Graph Context Protocol
(\gcp{}). \gcp{} is \emph{read-first}: an agent obtains what it needs by reading context from
independently owned knowledge nodes, rather than receiving pushed messages. It is
\emph{role-gated}: every read passes through a policy gate that authorizes the request by the
caller's role and required capabilities, defaulting to deny~\citep{saltzer1975,sandhu1996}. And
it is \emph{audited}: every read decision---grant or deny---produces a structured provenance
record naming the principal, target node, decision, and matched policy.

\paragraph{Thesis.} At equal task success, read-first role-gated context federation beats
point-to-point message-passing on coupling, leakage, and cost.

\paragraph{Contributions.}
\begin{enumerate}
  \item \textbf{Architecture.} We describe \gcp{}: read-first, role-gated context federation with
  first-class per-read provenance (Section~\ref{sec:arch}), and contrast it with push-based
  message-passing.
  \item \textbf{A fairness-controlled benchmark method.} The single most common flaw in
  architecture comparisons is that the arms differ in more than the architecture. Our harness
  builds both arms from one scenario definition and one injected model, swapping \emph{only} the
  interop layer (Section~\ref{sec:method}).
  \item \textbf{An empirical head-to-head} over three scenarios with a real LLM and five seeds,
  reporting counts (connections, messages, tokens), not just wall-clock (Section~\ref{sec:results}).
  \item \textbf{Metric definitions} for structural coupling, canary exact-match leakage (a lower
  bound), and provenance completeness (Section~\ref{sec:metrics}).
\end{enumerate}

We are deliberately conservative about the empirical claims: the behavioral results come from a
single free-tier model over five seeds and are best read as a \emph{pilot}. Section~\ref{sec:threats}
states the threats to validity in full, including a case where the architectural leak in \gcp{}'s
favor is \emph{latent rather than observed} because the model never exercised the leaky path.

\section{Background and Related Work}
\label{sec:related}

\subsection{Message-passing as the multi-agent baseline}
Agent communication languages established the message-passing paradigm: KQML~\citep{finin1994}
wrapped knowledge payloads in typed performatives exchanged over pairwise connections, and
FIPA-ACL~\citep{fipa2002} standardized a 12-field message in which sender and receiver must share
content language and ontology before any exchange succeeds. Modern LLM frameworks reinvented this
in natural language. AutoGen~\citep{wu2023autogen} programs applications as conversations between
agents; ChatDev~\citep{qian2024chatdev} arranges role-pairs in a fixed ``chat chain''; and surveys
of the field~\citep{guo2024survey} taxonomize communication structure into peer-to-peer,
centralized, layered, and shared-pool variants. The Model Context Protocol~\citep{mcp2024}
standardizes agent-to-tool context access (an early read-first form), while A2A~\citep{a2a2025}
standardizes agent-to-agent task delegation---the baseline we benchmark against. Recent work has
begun to quantify message-passing's overhead: \citet{zhang2024agentprune} report 28--73\% excess
tokens in inter-agent communication graphs and recover most of it by pruning the message topology.
Most directly relevant, \citet{mao2026delm} replace both centralized orchestration and peer-to-peer
routing with a shared verified context and report state-of-the-art results at roughly half the
cost---evidence that a read-from-shared-context substrate can dominate message-passing. We add to
that line role-gated authorization, audited per-read provenance, and a fairness-controlled head-to-head.

\subsection{Coupling and the $n^2$ problem}
That fewer inter-component links yield more maintainable systems is among the oldest results in
software engineering. \citet{parnas1972} showed that information-hiding decomposition---exposing
only a minimal interface---produces dramatically lower coupling than flowchart-style partitioning.
\citet{stevens1974} introduced coupling and cohesion as measurable, ordinal properties of
inter-module dependency. \citet{briand1999} gave a graph-theoretic, measurement-theoretic framework
in which coupling is a monotone function of dependency-edge count, so an $\bigoh(n)$-edge topology
is provably less coupled than an $\bigoh(n^2)$ one. \citet{hohpe2003} translated this to distributed
integration: a full point-to-point mesh of $n$ systems needs $n(n{-}1)/2$ channels, which a
hub-and-spoke broker reduces to $n$. Microservice research~\citep{panichella2021} carries the same
edge-count argument into cloud-native systems, and in the LLM setting \citet{shen2025topology} show
empirically that dense, fully-connected agent topologies amplify error propagation while sparser
topologies suppress it. We apply this directly: \ata{} creates pairwise
agent$\leftrightarrow$agent edges; \gcp{} creates agent$\rightarrow$store edges.

\subsection{Least privilege, RBAC, and gated context}
\gcp{}'s read gate operationalizes the principle of least privilege, articulated by
\citet{saltzer1975} and equated there with military need-to-know. Role-based access
control~\citep{sandhu1996} decouples permissions from individuals by assigning them to roles, and
attribute-based access control~\citep{nist800162} generalizes the policy to arbitrary attributes,
and capability-based protection~\citep{dennis1966} ties authority to unforgeable, explicitly passed
tokens---the model \gcp{}'s role credentials follow. Data-minimization and purpose-limitation
obligations~\citep{gdpr2016} make ``read only what your role requires'' a legal constraint, not
merely a design preference. The point-to-point baseline has
no analogous gate: leakage prevention rests entirely on sender-side discretion.

\subsection{Leakage, provenance, and governance in agentic systems}
The security literature establishes that confidentiality in agentic systems is a system property,
not only a model property~\citep{evertz2024whispers}, that inter-agent message channels are the
dominant leakage vector---leaking far more than final outputs and largely invisible to output-only
auditing~\citep{elyagoubi2026agentleak}---and that restricting what an agent can
\emph{see}---rather than filtering what it can \emph{say}---is the more robust
intervention~\citep{bagdasarian2024airgap}. The OWASP Top 10 for LLM
applications~\citep{owasp2025} names Sensitive Information Disclosure and Excessive Agency as
primary risks, with ``embrace least privilege'' as the normative mitigation. On the accountability
side, deterministic privilege control at the tool boundary~\citep{shi2025progent}, authenticated and
auditable delegation of authority to agents~\citep{south2025delegation}, and audit trails
for LLMs~\citep{ojewale2025audit} argue for tamper-evident, queryable records of what a system did.
\gcp{}'s contribution to this thread is to make the \emph{per-read decision} itself the provenance
primitive: every authorized or denied context read is recorded, which message-passing does not do.

\paragraph{Delta.} Prior work supplies the pieces---read-first access (MCP), role/attribute gating
(RBAC/ABAC, gated RAG), and audit trails---but not their conjunction, and not a fairness-controlled
head-to-head against a real message-passing baseline. \gcp{} combines read-first access, role
gating, and audited per-read provenance, and we measure the combination against \ata{}.

\section{The Graph Context Protocol}
\label{sec:arch}

\gcp{} models a multi-agent ecosystem as independently owned \emph{knowledge nodes} and
\emph{agent nodes}. Coordination is read-first: an agent acquires context by issuing a
\emph{context query} to a peer node, which returns context the caller is authorized to read. Three
properties distinguish the substrate.

\paragraph{Read-first.} Producers do not push to consumers. A knowledge node exposes context;
consumers pull on demand. No agent needs to know the identity of its producers' \emph{consumers},
and adding a consumer requires no change at the producer. Structurally, each agent connects to the
nodes it reads, not to every other agent.

\paragraph{Role-gated.} Each knowledge node carries an access policy
(\texttt{readableByRoles}, \texttt{requiredCapabilities}, \texttt{fallbackAllowed},
\texttt{denialMode}) and the server authorizes every query against the caller's principal, defaulting
to deny when no policy grants access. Authorization happens \emph{before} any content reaches the
agent's context window---the read-first analogue of complete mediation~\citep{saltzer1975}.

\paragraph{Audited.} Every read decision emits a structured provenance record
(\texttt{principalId}, \texttt{targetNodeId}, \texttt{timestamp}, \texttt{decision},
\texttt{matchedPolicy}) to a pluggable audit sink. This yields a complete who-read-what-when ledger
for the federation, including denials.

The message-passing baseline (\ata{}) implements the same scenario tasks by having agents hold
peer references via cards and exchange context through direct messages, with coarse card-declared
allow/deny and no per-read audit sink---a faithful, good-faith control built with the same care as
the \gcp{} arm.

\section{Threat and Cost Model, and Metric Definitions}
\label{sec:metrics}

We report metrics in three families.

\paragraph{Structural coupling (deterministic).} For $n$ peer nodes, the \emph{pairwise connections}
are the number of integration edges the topology requires. The point-to-point arm needs
$C_{\ata{}}(n)=n(n{-}1)$ directed edges (every agent$\leftrightarrow$peer pair); the federated arm
needs $C_{\gcp{}}(n)=n$ (each agent$\rightarrow$store). \emph{Integration effort} is the marginal
cost of adding one node, $C(n{+}1)-C(n)$: constant $1$ for \gcp{} and $2n$ for \ata{}. These are
properties of the topology, not of any run, so they are computed exactly.

\paragraph{Leakage (a lower bound).} Each confidential node embeds a unique \emph{canary} string.
After a run we scan the final answer \emph{and} the full tool transcript for exact canary matches;
\emph{leakage rate} is the fraction of a scenario's canaries that appear where they should not.
Exact-match scanning is a lower bound: paraphrased disclosure is not counted.

\paragraph{Provenance completeness.} The fraction of read decisions for which a structured
provenance record exists, clamped to $[0,1]$. \gcp{} records every decision; the message-passing
arm has no audit sink, so its completeness is $0$ by construction.

\paragraph{Behavioral cost and success.} We count messages, connections opened, round-trips, and
total tokens (summed from provider usage metadata), measure wall-clock latency, and evaluate task
success with a per-scenario oracle. Counts---not just wall-clock---are the primary cost signal,
since wall-clock on a shared free-tier endpoint is noisy.

\section{Methodology}
\label{sec:method}

\paragraph{The fairness invariant.} The central methodological commitment is that the two arms
differ \emph{only} in the interop layer. A single scenario runner builds both arms from one
\texttt{ScenarioDef} and one injected chat model. Both arms run the identical agent brain
(a ReAct agent with one per-peer tool), the identical system prompt, goal, and success oracle. The
\gcp{} arm hosts knowledge nodes in-process behind the real server's fetch handler with a static-token
auth provider and an in-memory audit sink; the \ata{} arm hosts the same nodes as real localhost A2A
servers. Only the per-peer tool factory and node hosting change. This removes the confound that makes
most architecture comparisons uninterpretable.

\paragraph{Scenarios.} \emph{software-org} (leakage): a contractor-role agent must summarize public
information about a project; a confidential engineering node carrying a canary is readable only by the
engineering role. \emph{supply-chain} (cross-owner): a logistics agent must state the shareable
shipping relationship between two independently owned nodes (a supplier and a manufacturer), each
holding an owner-scoped canary. \emph{marketplace} (coupling): $n$ seller nodes with descending
prices; the agent must identify the cheapest. The marketplace is parameterized by $n$ and drives the
structural-coupling curve.

\paragraph{Model, seeds, harness.} All behavioral runs use \texttt{openai/gpt-oss-120b:free} via
OpenRouter at temperature $0.7$, with five repetitions (seeds) per arm per scenario. The harness is
gated behind an environment flag and an API key so it never runs in CI; it throttles between runs and
retries on transient rate-limit errors, both of which are common on free-tier endpoints. We report
mean $\pm$ population standard deviation. The structural-coupling curve is computed deterministically
for $n=2,\dots,50$.

\section{Results}
\label{sec:results}

\subsection{Coupling scales as $\bigoh(n)$ vs.\ $\bigoh(n^2)$}
Table~\ref{tab:struct} shows the structural-coupling curve. \gcp{}'s pairwise connections grow
linearly ($n$) while \ata{}'s grow quadratically ($n(n{-}1)$): identical at $n{=}2$, $3\times$ at
$n{=}3$, $9\times$ at $n{=}10$, and $49\times$ at $n{=}50$ ($50$ vs.\ $2450$). The marginal cost of
adding a node---integration effort---is constant $1$ for \gcp{} and $2n$ for \ata{}, so the gap
\emph{widens} with every node added.

\begin{table}[t]
\centering
\caption{Structural coupling (pairwise integration edges) by ecosystem size $n$. Deterministic.}
\label{tab:struct}
\begin{tabular}{rrrr}
\toprule
$n$ & \gcp{} $(n)$ & \ata{} $(n(n{-}1))$ & ratio \ata{}/\gcp{} \\
\midrule
2  & 2  & 2    & 1.0$\times$ \\
3  & 3  & 6    & 2.0$\times$ \\
5  & 5  & 20   & 4.0$\times$ \\
10 & 10 & 90   & 9.0$\times$ \\
20 & 20 & 380  & 19.0$\times$ \\
50 & 50 & 2450 & 49.0$\times$ \\
\bottomrule
\end{tabular}
\end{table}

\subsection{Provenance is complete for \gcp{} and absent for \ata{}}
Across all three scenarios and all five seeds, provenance completeness is $1.0$ for \gcp{} and $0.0$
for \ata{} (Table~\ref{tab:behav}). This is architectural, not statistical: \gcp{} records a decision
for every read; the message-passing arm has no audit sink, so no access trail exists to query. For any
governance regime that requires traceability~\citep{ojewale2025audit,owasp2025}, this is the sharpest
qualitative difference between the substrates.

\subsection{\gcp{} never costs more, and costs less on the cross-owner task}
Table~\ref{tab:behav} reports behavioral metrics. On \emph{software-org} the two arms are within noise
on tokens ($678.4{\pm}5.8$ vs.\ $685.4{\pm}3.4$). On \emph{supply-chain} \gcp{} uses fewer tokens with
far lower variance ($707.6{\pm}3.5$ vs.\ $901.4{\pm}229.2$) and fewer messages/round-trips
($1.0{\pm}0.0$ vs.\ $1.4{\pm}0.5$): the message-passing arm sometimes opened extra connections to
assemble the cross-owner answer, while the federated arm read what it needed in one pass. Task success
is $1.0$ on both arms for these two scenarios---the equal-success precondition of the thesis holds.

\begin{table*}[t]
\centering
\caption{Behavioral metrics, mean $\pm$ population s.d.\ over 5 seeds, \texttt{gpt-oss-120b:free}.
Structural columns are for the scenario's node count ($n{=}2$).}
\label{tab:behav}
\setlength{\tabcolsep}{4pt}
\begin{tabular}{llrrrrr}
\toprule
scenario & arm & tokens & messages & latency (ms) & leakage & provenance \\
\midrule
\multirow{2}{*}{software-org}
 & \gcp{} & $678.4{\pm}5.8$  & $1.0{\pm}0.0$ & $6268{\pm}482$  & $0.0$ & $1.0$ \\
 & \ata{} & $685.4{\pm}3.4$  & $1.0{\pm}0.0$ & $5986{\pm}1168$ & $0.0$ & $0.0$ \\
\midrule
\multirow{2}{*}{supply-chain}
 & \gcp{} & $707.6{\pm}3.5$   & $1.0{\pm}0.0$ & $6278{\pm}1632$ & $0.5$ & $1.0$ \\
 & \ata{} & $901.4{\pm}229.2$ & $1.4{\pm}0.5$ & $7718{\pm}2678$ & $1.0$ & $0.0$ \\
\midrule
\multirow{2}{*}{marketplace$^\dagger$}
 & \gcp{} & $2425{\pm}1001$ & $2.8{\pm}0.7$ & $12617{\pm}3757$ & $0.0$ & $1.0$ \\
 & \ata{} & $2431{\pm}1082$ & $2.8{\pm}0.7$ & $11941{\pm}4304$ & $0.0$ & $0.0$ \\
\bottomrule
\end{tabular}
\par\smallskip
\footnotesize $^\dagger$ marketplace rows aggregate $n\in\{2,5,10\}$; success was $0.3{\pm}0.5$
on both arms (see Section~\ref{sec:threats}). Use the structural curve (Table~\ref{tab:struct}),
not these mixed-$n$ behavioral rows, for the coupling claim.
\end{table*}

\subsection{Leakage}
The leakage picture requires care. On \emph{supply-chain}, where the logistics role is authorized to
read both owners' nodes, the message-passing arm leaked every owner-scoped canary (rate $1.0$) while
\gcp{} leaked half on average ($0.5$); because both arms are authorized to read the content, this gap
reflects how each substrate surfaces content to the model rather than a gating decision. On
\emph{software-org}, the cleanest gating case, the \emph{observed} leakage was $0.0$ on both arms---but
for an instructive reason: with this model the agent satisfied the public-summary task by querying only
the public node (messages $=1.0$) and never requested the confidential node, so the leaky path was not
exercised. \gcp{} would have \emph{denied} that read by construction had it occurred (and recorded the
denial); the message-passing arm's coarse card would have answered it. The architectural leak is thus
real but \emph{latent} in this run, and we do not claim an observed behavioral leakage win on
software-org. The structural guarantee---deny-by-default plus audited denial---is what \gcp{} provides
regardless of whether a given model happens to probe the boundary. The canary metric \emph{measures}
this bound; the formal model and executable property suite described in
Section~\ref{sec:formal-no-leak} \emph{prove} it.

\subsection{Formal model and executable no-leak property}
\label{sec:formal-no-leak}

\paragraph{Access model.}
A knowledge node carries an \texttt{AccessPolicyDescriptor} $P$ with two gating fields:
$\textit{readableByRoles}$ (a list of role ids) and $\textit{requiredCapabilities}$ (a list of
capability ids). A principal $K$ holds an optional role with its own and inherited capability set
$\mathit{effective}(K)$, computed by \texttt{getEffectiveCapabilities} with own-first precedence when
the same capability id appears at multiple levels of the role chain.
The admission predicate is:
\[
\mathit{authorized}(P,K) \;\equiv\;
  \underbrace{\bigl(P.\textit{readableByRoles}=[\,] \;\vee\; K.\textit{role}.\textit{id}\in P.\textit{readableByRoles}\bigr)}_{\textit{roleOk}}
  \;\wedge\;
  \underbrace{\bigl(P.\textit{requiredCapabilities}=[\,] \;\vee\; P.\textit{requiredCapabilities}\subseteq\mathit{effective}(K)\bigr)}_{\textit{capsOk}}
\]
An absent or schema-invalid policy descriptor is treated as a denial (default-deny): the gate is
closed whenever a valid policy is not available.

\paragraph{No-leak theorem (informal).}
For all $P$ and $K$: if $\neg\mathit{authorized}(P,K)$ then no \texttt{context-query} directed at a
node carrying $P$ returns that node's content to $K$. The argument follows from the handler pipeline
in \texttt{createContextQueryHandler}: authorization (step~4) precedes any knowledge-source adapter
call (step~5). On failure the handler returns a \texttt{denied} result containing only the denial
reason; the adapter is never invoked and the node's content never enters the response.

\paragraph{Executable property suite.}
We implement three \texttt{fast-check} property-based tests over the \emph{real} shipped code path
at CI defaults of $\mathit{FC\_NUM\_RUNS}=1000$; all pass.

\textbf{P1~(end-to-end no-content-leak).}
For each randomly generated $(P,K)$ pair, the property wires a live
\texttt{createContextQueryHandler()} instance with a stub adapter that unconditionally returns a
sentinel string \texttt{SECRET-CANARY-\ldots} in its \texttt{raw} field and an allow-all auth
provider (so the provider term drops out). Two directions are asserted simultaneously:
\emph{safety} --- when $\neg\mathit{authorized}(P,K)$, the full JSON serialisation of the handler
response does not contain the sentinel and \texttt{payload.status} is \texttt{"denied"}; and
\emph{liveness} --- when $\mathit{authorized}(P,K)$, the sentinel \emph{does} appear, confirming the
harness is capable of leaking and is not vacuously passing.

\textbf{P2~(gate soundness).}
For each randomly generated $(P,K)$ pair, the real \texttt{authorizeKnowledgeNodeAccess} function
is called against a node carrying $P$ with an allow-all provider, and the result is asserted to equal
$\mathit{authorized}(P,K)$ computed by the independent reference oracle. This establishes that the
reference predicate is an accurate specification of the shipped gate, the key compositional link that
lets P1 use $\mathit{authorized}$ as its oracle.

\textbf{P3~(role-inheritance soundness).}
Random chains of roles (depth 2--4) over a fixed pool of five capability ids are generated and the
completeness and own-first-precedence invariants of \texttt{getEffectiveCapabilities} are verified.
P3 grounds the effective-capability computation used by both the reference predicate and the real
gate.

\paragraph{Canary measurement vs.\ formal proof.}
The canary metric (Section~\ref{sec:method}) quantifies how often leakage occurs across
production-like scenarios; it is a \emph{measured} lower bound. The property suite here is a
\emph{proof-level} complement: it proves that the no-leak bound holds over the entire randomised
input space (subject to finite generator pools). The canary metric can detect a regression if the
gate is bypassed; the property suite explains why no leak should occur and gives a constructive
account of the invariant. Together they form a two-layer assurance: the property proves the gate is
sound; the canary confirms the gate is on the hot path in production-like scenarios.

\section{Discussion and Threats to Validity}
\label{sec:threats}

\paragraph{Single model, five seeds.} The behavioral results use one free-tier model
(\texttt{gpt-oss-120b}) over five seeds. We therefore present them as a pilot. The structural-coupling
and provenance-completeness results do \emph{not} depend on the model---the former is deterministic,
the latter architectural---so they are robust to this limitation; the token and latency results should
be replicated across models before being treated as general.

\paragraph{Latent vs.\ observed leakage.} As Section~\ref{sec:results} details, the software-org leak in
\gcp{}'s favor was not behaviorally observed because the model did not probe the confidential node.
This is a genuine limitation of behavioral leakage measurement: it depends on the agent exercising the
path. We mitigate by also reporting the architectural guarantee and the provenance of denials, and the
harness includes a forced-query hermetic test that exercises the leaky path directly.

\paragraph{Leakage is a lower bound.} Exact-match canary scanning misses paraphrase and partial
disclosure. Reported leakage rates therefore understate true exposure for both arms equally.

\paragraph{Marketplace task difficulty.} Marketplace task success was low ($0.3$) for \emph{both} arms:
\texttt{gpt-oss-120b} often failed to identify the cheapest seller. Because the failure is equal across
arms it does not threaten fairness, but it means the marketplace contributes through its deterministic
structural curve, not its behavioral rows.

\paragraph{Free-tier noise.} Wall-clock latency on a shared free endpoint is high-variance; we lead with
counts (messages, tokens) for cost and treat latency as secondary.

\paragraph{Construct validity of the baseline.} A weak baseline would make the comparison vacuous. The
\ata{} arm runs real servers, the same brain, and the same task, and was built with equal care; its
disadvantages (no read gate, no audit sink, pairwise topology) are intrinsic to message-passing, not
artifacts of a strawman.

\section{Limitations and Future Work}
The evaluation should be extended across multiple models and more seeds, with significance testing on
the token and latency deltas. One designed-but-unevaluated capability remains: a stricter-capability
\emph{task-delegation} layer on the same read-first path. The formal model with an executable no-leak
property is now complete (Section~\ref{sec:formal-no-leak}), upgrading the leakage argument from
measured to proven. Bridging \gcp{} to MCP and A2A as first-class interop (beyond the benchmark
baseline) is the natural deployment path.

\section{Conclusion}
We argued, and gave preliminary fairness-controlled evidence, that read-first role-gated context
federation beats point-to-point message-passing for multi-agent LLM systems on coupling, leakage, and
cost at equal task success. Coupling improves from $\bigoh(n^2)$ to $\bigoh(n)$ with constant marginal
integration effort; provenance is complete by construction rather than absent; and cost is never worse
and sometimes better. The architecture turns leakage prevention from a matter of sender discretion into
a deny-by-default, audited read gate. The strongest results---coupling and provenance---are
model-independent; the cost and leakage results are a model-specific pilot that invites replication.

\bibliographystyle{unsrtnat}
\bibliography{refs}

\end{document}
